\begingroup
\begin{center}\textbf{Résumé}\end{center}
\fontsize{9.1}{10.55}\selectfont
\setlength{\parskip}{1.2pt}
\setlength{\abovedisplayskip}{3pt}
\setlength{\belowdisplayskip}{3pt}
On démontre l’hypothèse du continu dans la clôture généalogique construite
par l’holographie modulaire triadique (HMT) : chaque sous-ensemble interne de sa droite est
dénombrable ou a la cardinalité de la droite complète. HMT est le noyau formel
qui compose trois objets. All Possible Paths (APP) organise les évaluations
de somme et de produit sur un réseau modulaire, dont le graphe de Cayley est le
squelette d’une surface cellulaire toroïdale bidimensionnelle. TRIT est la signature
ternaire orientée \(\{-1,0,+1\}\). Le Topological Phase Kernel (TPK), ou noyau
topologique de phase, sélectionne et transporte les transitions, conserve leur
mémoire et prolonge les histoires compatibles. La limite inverse
\(X_\infty^{\mathrm{enr}}\) conserve valeur, feuillet, orientation, reste,
quotient, retenue, frontière, route, mémoire et survie. La structure
discrète conjointe du continu et sa droite archimédienne s’obtiennent comme sorties
de cette action.

Soit \(h\) le code effectif des graphes des opérateurs APP, TRIT et TPK, de leurs prolongements, de leurs restrictions et des deux projections. Pour chaque histoire réalisée \(m_\infty\), la clôture généalogique \(\mathfrak G_{\mathrm{HMT}}(h,m_\infty)\) s'obtient par itération transfinie de définitions finies avec des paramètres de niveaux antérieurs. Le théorème principal établit
\[
 \boxed{\mathfrak G_{\mathrm{HMT}}(h,m_\infty)
 \models 2^{\aleph_0}=\aleph_1}
\]
et, pour tout \(A\in\mathcal P^{\mathfrak G}
(\mathbb R_{\mathrm{HMT}}^{\mathfrak G})\), prouve la dichotomie
\[
 |A|\leq\aleph_0
 \qquad\text{ou}\qquad
 |A|=|\mathbb R_{\mathrm{HMT}}^{\mathfrak G}|.
\]
Par conséquent, il n'existe pas de cardinal intermédiaire entre les entiers naturels et la droite à l'intérieur du continu intégral engendré par HMT.

La démonstration enchaîne quatre étapes. La couverture exhaustive des cylindres établit l'équipotence de la droite obtenue par réduction de cylindres avec \(\mathcal P^{\mathfrak G}(\omega)\). La condensation généalogique construit une injection \(j_{\mathrm{HMT}}:\mathbb R_{\mathrm{HMT}}^{\mathfrak G}
\hookrightarrow\omega_1^{\mathfrak G}\). Le codage des bons ordres dénombrables et leur élévation ternaire construisent l'injection en sens contraire \(k_{\mathrm{HMT}}\). Cantor--Bernstein conclut l'égalité cardinale, et l'argument de cofinalité étend la conclusion à tous les sous-ensembles de la clôture complète. Aucune de ces étapes ne présuppose l'égalité qu'elle démontre.

La fibre locale de \(3^6=729\) sous-cylindres de raffinement ne borne pas le
résultat : la loi coinductive agit à toute profondeur prescrite. Le retour
nonadique retrouve la phase tandis que la retenue et la mémoire avancent ; il ne réinitialise pas
l’état. La même généalogie produit, comme résultats corrélés,
l’holonomie \(\Gamma_9\), le registre dodécaphasique entier \(K\), la
quadrirelation \((\pi,\varphi,e,\alpha)\) et la réalisation incidencielle
Paley--Witt--Golay--Leech--\(V^\natural\)--Monstre. Ces sorties ne prouvent pas
l’égalité cardinale par coïncidence numérique ; elles certifient l’unité de
l’état dont la projection archimédienne détermine la valeur et dont la projection
incidencielle conserve la structure.

Enfin,
\(\mathfrak G_{\mathrm{HMT}}(h,m_\infty)=L[h,m_\infty]\) est démontré comme
reconnaissance ultérieure de la hiérarchie déjà construite. Il n’engendre pas la droite,
ne sélectionne pas les injections et n’intervient pas dans la réduction des cylindres.

\medskip

% BEGIN R27_FRAMING_SUMMARY_ES
La description de l’émetteur est complétée par la preuve de son alphabet
décimal à \(42\) blocs, un cycle explicite d’incidences entre
semi-émissions régionales et une comparaison de mémoire ordonnée à
histogramme fixe. Ces conséquences spécifient l’information
conservée par les lectures finies qui précèdent les prolongements.
La comparaison cardinale maintient le domaine et les démonstrations
propres à la clôture généalogique.
% END R27_FRAMING_SUMMARY_ES

\noindent\textbf{Mots-clés :} holographie modulaire triadique ; APP ; TRIT ; TPK ; état enrichi ; structure discrète du continu ; clôture généalogique ; hypothèse du continu ; cardinalité ; ordinalité ; coinduction.
\endgroup
